\section{Chuyển đổi sang TensorRT}

\subsection{Export ONNX}

Mỗi checkpoint ứng viên được export sang ONNX với đầu vào tĩnh
$1\times3\times640\times640$ (batch 1, phù hợp xử lý luồng camera từng khung)
bằng chức năng export của Ultralytics. ONNX checker xác nhận đồ thị hợp lệ; suy
luận đối chiếu giữa PyTorch và ONNX Runtime trên một tập ảnh cố định đo sai khác
tensor đầu ra và sai khác detection cuối. YOLO26 có head end-to-end không cần NMS
nên đồ thị ONNX không chứa bước NMS phụ thuộc dữ liệu, giúp chuyển đổi TensorRT
thuận lợi.

\subsection{TensorRT FP16 và INT8}

Engine được build \emph{trên chính thiết bị đích} vì engine TensorRT gắn với
kiến trúc GPU và phiên bản TensorRT: engine tạo trên Jetson Nano (TensorRT 8.2)
không chạy trên Jetson Orin Nano và ngược lại. Engine FP16 được tạo từ baseline và
các mô hình đã cắt tỉa. Engine INT8 được tạo từ ONNX Q/DQ của mô hình QAT; vì
thang lượng tử đã có trong đồ thị nên không cần tập hiệu chỉnh bổ sung. Log build
ghi lại phiên bản TensorRT, workspace và precision thực tế của từng lớp để xác định
tỷ lệ phép toán thực sự chạy INT8.

\subsection{Kiểm tra độ chính xác sau chuyển đổi}

Kiểm tra gồm ba mức: sai khác tensor trên mẫu cố định, sai khác detection theo
IoU và confidence, và chênh lệch $mAP$ trên toàn tập test. Artifact vi phạm ràng
buộc độ chính xác ở Mục~3.1 bị loại dù latency tốt. Engine đạt yêu cầu được đẩy
lên kho model và triển khai qua nền tảng; AI Runtime nạp trực tiếp tệp
\texttt{.engine} bằng Ultralytics.
